% =========================================================
% frontmatter/abstract.tex
% =========================================================

\chapter*{Abstract}
\thispagestyle{fancy}
\addcontentsline{toc}{chapter}{Abstract}
\begingroup
\setlength{\emergencystretch}{2em}

This report presents the analytical design and CST full-wave simulation of a
finite-radius, center-fed wire dipole operating near 1~GHz. A first-order free-space
calculation gives a wavelength of 300~mm and an initial half-wave length of 150~mm.
The CST model uses a 5~mm-diameter PEC conductor, a 20~mm center gap, a 73~$\Omega$
discrete port, and open-add-space boundaries. Three documented lengths---150~mm,
125~mm, and 135.6~mm---establish the tuning path.

The initial 150~mm model resonated below the 1~GHz target, while shortening the
dipole to 125~mm shifted the resonance above the target. The final 135.6~mm model
produced a CST marker at 0.99648~GHz with
$S_{11,\min}=-47.01186$~dB, corresponding to a positive return loss of
47.01186~dB relative to the 73~$\Omega$ port reference. The frequency deviation from
1~GHz is 3.52~MHz, or 0.352\%.

The final CST far-field result displays the expected broadside toroidal pattern with
axial nulls and a peak directivity of 2.149~dBi. Independent evaluation of the canonical
thin half-wave dipole pattern gives 2.15~dBi directivity and a 78.08$^\circ$ half-power
beamwidth, in close agreement with the displayed CST pattern. The study therefore
connects closed-form antenna theory, parameterized full-wave modeling, resonance
control through length tuning, and analytical-to-simulation verification.

All performance results are analytical or simulated. No fabricated prototype or
measured antenna data are claimed. The absence of the native CST model, raw exports,
and formal mesh- and boundary-convergence records limits the conclusions to the
preserved evidence package.
\par\endgroup
